\documentclass[10pt,a4paper]{amsart}
\usepackage[margin=19mm]{geometry}
\usepackage{amsmath,amssymb,mathtools}
\usepackage[scr=rsfs,cal=euler]{mathalfa}
\usepackage{xfrac}
\usepackage[unicode,hidelinks,bookmarks=false]{hyperref}
\newcommand{\ie}{i.e.\ }
\newcommand{\cf}{cf.\ }
\newcommand{\resp}{respectively\ }
\newcommand{\Subsection}{\subsection}
\providecommand{\nobreakdash}{\nobreak\hbox{-}}
\setlength{\emergencystretch}{2em}
\multlinegap0pt
\usepackage{tikz}
\usepackage{tikz-cd}
\usepackage{latexsym}
\usepackage{graphics}
\usepackage{fancyhdr}
\usepackage{epstopdf}
\usepackage{setspace}
\usepackage{url}
\usepackage{youngtab}
\newtheorem{theorem}{Theorem}[section]
\newtheorem{cor}[theorem]{Corollary}
\newtheorem{prop}[theorem]{Proposition}
\newtheorem{lemma}[theorem]{Lemma}
\newtheorem{claim}[theorem]{Claim}
\newtheorem{conj}[theorem]{Conjecture}
\newtheorem{corollary}[theorem]{Corollary}
\newtheorem{definition}[theorem]{Definition}
\newtheorem{question}[theorem]{Question}
\def\blkbx{\raisebox{-1mm}{\rule{2mm}{4mm}}}
\newenvironment{dfn}{\refstepcounter{theorem} \bigskip \par \noindent
                    {\bf Definition \thetheorem}\ }{\bigskip \par}
\def\R{{\mathbb R}} % Reals
\def\C{{\mathbb C}} % Complex
\def\P{{\mathbb P}} % Projective Space
\def\A{{\mathbb A}} % Affine Space
\def\Z{{\mathbb Z}} % Integers
\def\N{{\mathbb N}} % natural numbers
\def\Q{{\mathbb Q}} % rational numbers
\begin{document}
\title[The Natural Extension for the Triangle Map (a Multi-dimensio]{The Natural Extension for the Triangle Map (a Multi-dimensional Continued Fraction) with An Internal Symmetry from Young Conjugation}
\author{Joe Fox and Thomas Garrity; Williamstown, MA 01267; Joe Fox; Thomas Garrity}
\date{}
\maketitle
\noindent\emph{Source and licence.} The Natural Extension for the Triangle Map (a Multi-dimensional Continued Fraction) with An Internal Symmetry from Young Conjugation. Version arXiv:2606.28014v2 (CC BY 4.0). This material is distributed under the Creative Commons Attribution 4.0 International licence, http://creativecommons.org/licenses/by/4.0/, as recorded in the arXiv OAI licence metadata for this exact version and shown on the abstract page https://arxiv.org/abs/2606.28014v2. Full rights notice: licenses/arxiv-2606.28014-CC-BY-4.0.txt.\par\smallskip
\noindent\textit{Editorial scope.} Counted unit: the original \texttt{theorem} environment \texttt{-} at source lines 2203--2206 together with its complete proof at lines 2207--2208; the delivered document keeps source lines 2160--2209 so that every referenced label and every citation resolves inside it. Native editable source is the paper's own concatenated TeX; the AMS wrapper is editorial formatting only. No publisher class, upstream script or unrelated artwork is redistributed.
\section{The invariant measure} \label{The invariant measure} 

\subsection{The relevant Jacobians}
 Consider
the simplex 
$$\triangle_m^A= \{ (t_1, \ldots , t_m) \in \R^m:  1>t_1 > \cdots > t_m >0\}.$$
We have the two maps
$$T:\triangle_m^A \rightarrow \triangle_m^A$$
and
$$G:\triangle_m^A \rightarrow \triangle_m^A$$
defined by 

$$T(t_1, \ldots, t_m) = \left\{  \begin{array}{ccc} \left( \frac{t_2}{t_1} , \ldots, \frac{t_m}{t_1}, \frac{1-t_1}{t_1} \right)& \mbox{if}& 1< t_1 + t_m  \\
 \left( \frac{t_1}{ 1 - t_m} , \ldots, \frac{t_m}{1- t_m} \right)& \mbox{if}& 1> t_1 + t_m \end{array} \right.$$
 and 
 $$G(t_1, \ldots, t_m) =   \begin{array}{ccc} \left( \frac{t_2}{t_1} , \ldots, \frac{t_m}{t_1}, \frac{1-t_1 - k t_m}{t_1} \right)& \mbox{if}& 1-t_1-kt_m \geq 0 >1-t_1-(k+1)t_m   \end{array} .$$
 
 
\begin{prop} We have for all points $\overline{t} = (t_1, \ldots, t_m) \in \triangle_m^A$ that 
$$\mbox{Jacobian}(G) (\overline{t}) = \frac{1}{t_1^{m+1}}.$$
If $ 1< t_1+ t_m$, then 
$$\mbox{Jacobian}(T) (\overline{t}) = \frac{1}{t_1^{m+1}},$$
and if $1>t_1+ t_m$, then
$$\mbox{Jacobian}(T) (\overline{t}) = \frac{1}{(1-t_m)^{m+1}},$$
\end{prop}
\begin{proof}
These are calculations.
\end{proof}
 \begin{corollary} Let $\mu$ be Lebesgue measure on the simplex $\triangle_m^A$.  Then for measurable sets $U$ in $\triangle_m^A$, it is not necessarily the case that 
 $$\mu(U) = \mu(T^{-1}(U))$$
 or that 
 $$\mu(U) = \mu(G^{-1}(U)).$$
Thus the Lebesgue measure is not an invariant measure for either  map $T$ or for map $G$.
 \end{corollary}
 
The natural question is to find an invariant measure.  This is classical for $m=1$, done in \cite{Ass} for the Gauss map, and completely worked out in \cite{Garrity-Lehmann Duke}, where it was shown that in Theorem 7.1 that the invariant measure for $G$ is 
$$\frac{\mathrm{\mu}}{t_1 t_2 \ldots t_{m-1} ( 1 + t_m)}$$
(perversely, in that paper the map $G$ was denoted by $T$) and in Theorem 8.3 the invariant measure for our $T$ is 
$$\frac{\mathrm{\mu}}{t_1 t_2 \ldots t_{m}}.$$
The method in \cite{Garrity-Lehmann Duke}  for deriving these invariant measures was the time-honored approach of guessing, then checking.

There is another approach for finding invariant measures if the natural extension is known.  This is highlighted in Arnoux and Nogueira \cite{Arnoux-Nogueira-93}.  We start with observing that 


\begin{theorem} For all of our  natural extension maps, we have
$$|\mbox{Jacobian}(\tilde{T}_0^A )| = |\mbox{Jacobian}(\tilde{T}_1^A )| =  |\mbox{Jacobian}(\tilde{G}_n^A )| =1. $$

\end{theorem}

\begin{proof} These are calculations.
\end{proof}


\begin{thebibliography}{99}
\bibitem[Arnoux-Nogueira-93]{Arnoux-Nogueira-93} P. Arnoux and A. Nogueira, Mesures de Gauss pour des algorithmes de fractions continues multidimensionnelles, \emph{Ann. Sci. \'Ecole Norm. Sup. (4)}, \textbf{26} (1993), no. 6, 645--664.
\bibitem[Ass]{Ass} Assaf, Sami, Li-Chung Chen, Tegan Cheslack-Postava, Benjamin Cooper, Alexander Diesl, Thomas Garrity, Mathew Lepinski, and Adam Schuyler. A Dual Approach to Triangle Sequences: A Multidimensional Continued Fraction Algorithm. Electronic Journal of Combinatorial Number Theory, 5 (2005), A08. www.emis.de/journals/INTEGERS/papers/f8/f8.pdf.
\bibitem[Garrity-Lehmann Duke]{Garrity-Lehmann Duke} T. Garrity and J. Lehmann Duke, Ergodicity and algebraicity of the fast and slow triangle maps, {\it Ergodic Theory and Dynamical Systems} Vol. 46 (2026), no. 1, pp. 93-127.
\end{thebibliography}
\end{document}
